\documentclass[12pt, letterpaper]{article}

% Shared preamble (palette, macros, geometry)
\input{../../../Preambles/header.tex}

% Document-level: natbib (job-market-paper style)
\usepackage[round, sort&compress, authoryear]{natbib}
\bibliographystyle{aer}  % or chicago, apalike, plainnat — adjust per journal

% Title block
% \\title{Property Tax Assessment Regressivity\\thanks{Saani Rawat, Marquette University. Email: saani.rawat@marquette.edu. I thank \\dots\\ for helpful comments. All errors are my own.}}
% \\author{Saani Rawat\\Marquette University}
\title{Property Tax Assessment Regressivity}
% \author{Anonymous Author for Review}
\date{\today\\\textit{Preliminary and incomplete --- please do not cite without permission}}

\begin{document}
\maketitle

\begin{abstract}
\noindent Property-tax assessments are systematically regressive: within a taxing jurisdiction, lower-value homes are assessed at a higher fraction of market value than higher-value homes. Using 11.9 million arms-length residential sales from national CoreLogic micro-data, we replicate this fact---a within-county elasticity of the assessment ratio with respect to sale price of $-0.42$---and ask where in the assessment process it originates. Decomposing the regressivity spatially, we find it is a \emph{within-neighborhood} phenomenon: adding census-tract fixed effects \emph{steepens} the slope to $-0.52$, so the standard jurisdiction-level estimate understates the within-neighborhood inequity by roughly twenty percent. We rule out the most intuitive alternative---that low-value homes carry staler assessments because they trade less often---which a covariate decomposition shows accounts for under one percent of the regressivity. The pattern points to cross-sectional mass-appraisal error that anchors each parcel toward its neighborhood mean, rather than to temporal staleness or neighborhood-level mispricing. \\[4pt]
\textbf{JEL classification:} H71, H22, R51 \\
\textbf{Keywords:} property tax, assessment regressivity, vertical equity, mass appraisal, CoreLogic
\end{abstract}

\section{Introduction}\label{sec:intro}

The property tax is the single largest source of own-source revenue for U.S.\ local governments and the fiscal backbone of public education. Its legitimacy rests on a simple premise: that a property's assessed value tracks its market value, so that owners of equally valuable homes pay equal tax and the burden rises in proportion to value. A growing body of evidence shows this premise fails systematically. Within a single taxing jurisdiction, lower-value homes are assessed at a higher fraction of their market value than higher-value homes, so the effective tax rate \emph{falls} with price \citep{berry2021reassessing, mcmillen2020assessment}. This assessment regressivity is a regressive transfer among owners of the same jurisdiction; it is correlated with race \citep{avenancioleon2022assessment} and with the tax-foreclosure crises that have hollowed out cities such as Detroit \citep{atuahene2019taxed, hodge2017assessment}.

The magnitude of the regressivity is now well documented. \citet{berry2021reassessing} estimates a within-jurisdiction elasticity of the assessment ratio with respect to sale price of about $-0.37$ on a national sample; we replicate it at $-0.42$ on 11.9 million arms-length CoreLogic sales (Section~\ref{sec:results}). Yet the appropriate policy response depends on a question the literature has only begun to answer: \emph{where in the assessment process does the regressivity originate?} If assessed values are simply stale, the remedy is more frequent reassessment \citep{mikesell1980property, hou2023assessment}; if the valuation model is miscalibrated, the remedy is better mass appraisal; if the appeals process favors wealthier owners, the remedy is procedural \citep{mcmillen2013effect, doerner2014empirical}. Each diagnosis implies a different reform.

This paper uses national, parcel-level CoreLogic micro-data to localize the regressivity and adjudicate among mechanisms. Rather than asking only how large the regressivity is, we ask at what \emph{spatial scale} and through what \emph{channel} it operates. The closest prior decomposition, \citet{amornsiripanitch2022why}, partitions national regressivity into a valuation-methods component and a reappraisal-frequency component. We instead decompose it \emph{spatially}---within versus across neighborhoods---and directly test the most intuitive temporal channel.

Our central result is that assessment regressivity is a \emph{within-neighborhood} phenomenon, and that the standard jurisdiction-level estimate understates it. Adding census-tract fixed effects to the canonical within-county specification does not attenuate the price slope; it \emph{steepens} it, from $-0.41$ to $-0.52$ (Figure~\ref{fig:ladder}, Table~\ref{tab:decomp}). Comparing two homes on the same block, the cheaper home is over-assessed by even more than the county-wide figure implies. The intuition is that across neighborhoods, price differences partly reflect amenities that assessors do capitalize into value, which flattens the slope; within a neighborhood, the mass-appraisal model anchors every parcel toward a common neighborhood prediction, mechanically over-assessing the below-median home. The regressivity is thus a property of \emph{cross-sectional} valuation within neighborhoods, and conditioning on the neighborhood reveals roughly twenty percent more of it than the jurisdiction-level estimate.

We then rule out the most natural alternative: that low-value homes are over-assessed because they trade less often and therefore carry staler assessed values. Using the complete transaction history of each parcel, a covariate (mediation) decomposition shows that transaction frequency and recency together account for under one percent of the regressivity (Table~\ref{tab:h6}). The operative channel is not temporal staleness from infrequent trading; it is the valuation model's treatment of the price cross-section. This null sharpens the target for the mechanism analysis that follows.

These findings make three contributions. First, relative to \citet{berry2021reassessing} and the measurement literature \citep{mcmillen2023measures}, we localize the regressivity to the neighborhood and show that the jurisdiction-level estimate is a \emph{lower bound} on the within-neighborhood inequity. Second, relative to \citet{amornsiripanitch2022why}, we provide a spatial decomposition orthogonal to the valuation-versus-reappraisal split, and we directly reject the transaction-frequency channel. Third, our spatial result speaks to \citet{avenancioleon2022assessment}: they show the \emph{racial} assessment gap is roughly half a between-neighborhood phenomenon, whereas we find the \emph{price} regressivity is overwhelmingly within-neighborhood. Race operates between neighborhoods; price regressivity operates within them.

Two caveats bound the analysis. First, our assessment snapshot is of 2008--2009 vintage, so we restrict attention to sales in 2007--2010, where the recorded assessment is approximately contemporaneous with the sale; this window lies inside the housing bust, which may amplify measured regressivity and which we probe directly. Second, \citet{mcmillen2023measures} show that the standard regression-based regressivity measure carries a mechanical bias; we therefore confirm the within-neighborhood pattern with a distribution-based statistic. The remainder of the paper proceeds as follows. Section~\ref{sec:data} describes the data; Section~\ref{sec:strategy} the empirical strategy; Section~\ref{sec:results} presents the decomposition and the mechanism tests; Section~\ref{sec:discussion} discusses implications for assessment policy; and Section~\ref{sec:conclusion} concludes.

\section{Data}\label{sec:data}

Our analysis draws on CoreLogic's bulk academic data, licensed through the University of Cincinnati \citep{corelogic_data}. We use two record systems: the \emph{Owner Transfer} file, which records deed transactions (sale price, transfer date, deed and transaction type), and the \emph{Property Characteristics} file, which records each parcel's assessor data (assessed value, tax amount, assessment year, structural attributes, and geocoding). The two are linked at the parcel level by CoreLogic's persistent property identifier (CLIP).

\paragraph{Sample construction.} We restrict attention to arms-length residential sales. Following standard sales-ratio practice \citep{iaao2013ratio}, we drop intra-family transfers, foreclosure and REO sales, and transactions below \$5{,}000, which are unlikely to reflect market value. The dependent variable throughout is the \emph{assessment ratio}, the recorded assessed value divided by the sale price; following \citet{berry2021reassessing} we trim ratios outside $[0.01, 5.0]$ to remove data errors. Because the University of Cincinnati assessment snapshot is predominantly of 2008--2009 vintage, we restrict to sales transacted in 2007--2010, the window in which the recorded assessment is approximately contemporaneous with the sale; outside this window the ratio would conflate genuine assessment regressivity with appreciation drift between the assessment date and the sale. The resulting estimation sample contains 11.9~million arms-length sales across 1{,}773 county jurisdictions.

\paragraph{Geography.} For the spatial decomposition we map each parcel to its census tract using CoreLogic's \texttt{census\_id} field, which is populated for 92\% of the estimation sample and yields 53{,}903 distinct tracts across 45 states. County and tract identifiers nest, so census-tract fixed effects subsume county fixed effects.

\paragraph{Caveats.} Two features of the data bound the analysis and are addressed directly below. The 2007--2010 window lies inside the housing downturn, during which assessed values---updated infrequently by statute---lagged falling market prices; this may amplify measured regressivity, a possibility we probe in Section~\ref{sec:discussion}. And the assessment-vintage restriction means our estimates speak to the late 2000s; extending the panel into the 2010s would require a more current or historical-panel extract.

\section{Empirical Strategy}\label{sec:strategy}

\paragraph{The regressivity elasticity.} Our point of departure is the canonical within-jurisdiction specification \citep{berry2021reassessing}:
\begin{equation}
\log(\text{ratio}_{ij}) = \alpha_j + \beta \, \log(\text{price}_{ij}) + \varepsilon_{ij},
\label{eq:main}
\end{equation}
where $i$ indexes a parcel-transaction, $j$ a jurisdiction, $\text{ratio}_{ij}$ is the assessment ratio, $\text{price}_{ij}$ the sale price, and $\alpha_j$ a jurisdiction fixed effect. The coefficient $\beta$ is the within-jurisdiction elasticity of the assessment ratio with respect to value, and $\beta < 0$ is regressivity. Because the statutory tax rate is constant within a jurisdiction, it is absorbed by $\alpha_j$, so~\eqref{eq:main} delivers the same within-jurisdiction elasticity whether the outcome is the assessment ratio or the effective tax rate. We cluster standard errors by county throughout.

\paragraph{Spatial decomposition.} To localize the regressivity we re-estimate~\eqref{eq:main} varying only the level at which $\alpha_j$ is defined: no fixed effect (pooled), the county, and the census tract. Comparing the county and tract estimates on a common sample isolates the spatial scale at which regressivity operates. Let $\beta_{\text{county}}$ and $\beta_{\text{tract}}$ denote the two slopes. If regressivity arose because assessors misprice whole neighborhoods, conditioning on the tract would absorb it and $\beta_{\text{tract}}$ would shrink toward zero; if instead it operates within neighborhoods, $\beta_{\text{tract}}$ would be no smaller---and, as we show, is in fact larger---than $\beta_{\text{county}}$. We summarize the result with the across-tract share, $(\beta_{\text{county}} - \beta_{\text{tract}})/\beta_{\text{county}}$.

\paragraph{Mechanism test: transaction frequency.} To ask whether regressivity reflects temporal staleness, we add a candidate mediator $m_{ij}$ to~\eqref{eq:main}---either the years since the parcel's prior recorded sale or its lifetime transaction count---and measure how much the price slope changes. In the spirit of a Gelbach covariate decomposition, the mediated share $(\beta - \beta')/\beta$, where $\beta'$ is the slope with $m_{ij}$ included, is the fraction of the price--ratio relationship that operates through transaction frequency. This is a descriptive decomposition; we do not interpret $m_{ij}$ as causally manipulable.

\paragraph{Measurement robustness.} \citet{mcmillen2023measures} show that regression-based regressivity measures can be mechanically biased. We therefore corroborate the spatial pattern with a rank statistic that does not depend on the slope in~\eqref{eq:main}: the size-weighted mean, within county versus within tract, of the Spearman rank correlation between sale price and the assessment ratio. Because it uses only ranks, it is immune by construction to the mechanical bias of the regression measure.

\section{Results}\label{sec:results}

\subsection{Replication and the spatial decomposition}

Figure~\ref{fig:ladder} presents the central result. The elasticity of the assessment ratio with respect to sale price is small when estimated by pooled OLS ($-0.18$), rises to $-0.42$ once we absorb county fixed effects---replicating \citet{berry2021reassessing} to within $0.05$---and rises \emph{further}, to $-0.52$, once we absorb census-tract fixed effects. Because tract fixed effects strengthen rather than weaken the slope, the regressivity is not driven by assessors mispricing whole neighborhoods; it is sharpest \emph{within} neighborhoods. Table~\ref{tab:decomp} reports the estimates. Columns~(3) and~(4) are estimated on the identical tract-covered subsample, so the move from $-0.41$ to $-0.52$ is a clean within-neighborhood comparison: conditioning on the neighborhood uncovers roughly twenty percent more regressivity than the jurisdiction-level estimate.

\begin{figure}[h!]
  \centering
  \includegraphics[width=0.85\linewidth]{figures/fig1_decomposition.pdf}
  \caption{Within-jurisdiction price-regressivity by fixed-effects specification. Points are the estimated elasticity of the log assessment ratio with respect to log sale price; whiskers are 95\% confidence intervals (standard errors clustered by county). The slope steepens monotonically as the comparison set narrows from the full sample (pooled) to the county to the census tract.}
  \label{fig:ladder}
\end{figure}

\input{tables/table1_main_decomposition.tex}

\subsection{Ruling out transaction frequency}

A natural explanation for within-neighborhood regressivity is temporal: cheaper homes may transact less often and so carry staler assessed values that lag a rising market. We test this directly. Table~\ref{tab:h6} reports a covariate (mediation) decomposition of the price slope, adding either the years since the parcel's prior sale (a backward-looking staleness measure) or the parcel's lifetime transaction count. Neither mediator moves the base coefficient: the share of regressivity operating through transaction frequency is $0.0$--$0.5\%$. Temporal staleness from infrequent trading is not the channel.

\input{tables/table2_h6_mediation.tex}

\subsection{Robustness}

The within-neighborhood result is not an artifact of thinly-sampled tracts or of the regression measure. Restricting to the 47{,}688 census tracts with at least twenty sales leaves the tract-fixed-effects slope essentially unchanged ($-0.516$). And a distribution-free rank statistic tells the same story as the regression: the size-weighted mean within-group Spearman correlation between sale price and the assessment ratio is $-0.48$ within counties and $-0.53$ within tracts---more negative inside the neighborhood, and immune by construction to the mechanical bias of \citet{mcmillen2023measures}.

\section{Discussion}\label{sec:discussion}

\paragraph{What the pattern implies about the assessment process.} Taken together, the two findings point to a specific source of regressivity. It is not that low-value parcels carry staler assessments from infrequent trading---transaction frequency explains under one percent of the slope. Nor is it that assessors misprice whole neighborhoods---conditioning on the tract \emph{strengthens} the regressivity. What remains is the cross-sectional behavior of the valuation model \emph{within} a neighborhood: mass-appraisal systems predict each parcel's value from neighborhood-level comparables and characteristics, compressing predictions toward a common neighborhood value and so over-assessing the below-median home and under-assessing the above-median one. In this account, regressivity is a feature of how automated valuation handles within-neighborhood heterogeneity, not of reassessment timing.

\paragraph{Implications for assessment policy.} This diagnosis bears on the remedy. The most-discussed reform---more frequent reassessment \citep{mikesell1980property, hou2023assessment}---targets temporal staleness, the channel we find to be quantitatively negligible. The pattern instead implicates the valuation model itself: its functional form, its delineation of neighborhoods, and its handling of within-neighborhood price dispersion. It also implies that the jurisdiction-level diagnostics on which assessors and regulators rely \citep{iaao2013ratio} understate the inequity a homeowner experiences relative to comparable neighbors; a within-neighborhood standard reveals roughly a fifth more regressivity than the county-level figure.

\paragraph{Relation to racial inequality in taxation.} \citet{avenancioleon2022assessment} show that the racial assessment gap is roughly half a between-neighborhood phenomenon. Our result is its complement on the price dimension: the price regressivity is overwhelmingly within-neighborhood. The two are consistent and jointly informative. Across neighborhoods, assessments track the amenities and composition that sorting capitalizes into prices; within neighborhoods, they fail to track the parcel-level price variation that remains. A complete account of who is over-taxed needs both margins.

\paragraph{Caveats.} Our estimates come from sales in 2007--2010, a window inside the housing bust. A downturn can inflate \emph{level} measures of regressivity, because assessed values lag falling prices; but our central claim concerns the cross-section \emph{within neighborhoods}, which the rank statistic confirms is not an artifact of the regression and which is less sensitive to the aggregate price path than a level comparison. Extending the analysis into the post-2010 recovery---which requires a more current or historical-panel assessment extract---would establish whether the within-neighborhood pattern is a crisis-era feature or a structural one. Finally, the present analysis is reduced-form; a structural model of the assessor's valuation problem, disciplined by these facts, would quantify the revenue misallocation by price decile and trace it to counterfactual reforms.

\section{Conclusion}\label{sec:conclusion}

Property-tax assessment regressivity is large, national, and---we have argued---fundamentally local. Replicating the canonical within-jurisdiction estimate on 11.9 million CoreLogic sales, we show that the regressivity is sharpest \emph{within} census tracts: conditioning on the neighborhood raises the elasticity of the assessment ratio with respect to price from $-0.41$ to $-0.52$, so the standard jurisdiction-level figure is a lower bound on the inequity a homeowner faces relative to comparable neighbors. The pattern is not driven by infrequent trading, which accounts for under one percent of the slope, and it survives a distribution-free rank statistic immune to the mechanical bias of regression-based measures. The evidence points to within-neighborhood mass-appraisal error---not temporal staleness or neighborhood-level mispricing---as the operative channel, reframing the policy problem from \emph{when} assessors reassess to \emph{how} their valuation models treat the price cross-section within a neighborhood.

\bibliography{../../../Bibliography_base}

\input{appendix.tex}

\end{document}
